\documentclass[a4paper]{article}

\usepackage{graphicx}
\usepackage{amsmath,amssymb}
\usepackage{minted}
\usepackage{multirow}
\usepackage{float}
\usepackage{todonotes}
\usepackage{forest}
\usepackage{xstring}
\usepackage{xcolor}
\usepackage[most]{tcolorbox}
\usepackage{xparse}
\usepackage{natbib}

\usetikzlibrary{graphs,graphdrawing}
\usegdlibrary{layered}

% for external graphviz
\input{/home/donkarlo/Dropbox/repo/nd_language_ai_project/src/nd_language_ai/formal/computer/programming/tex/latex/code_clips/external_graphviz.tex}

% for tags
\input{/home/donkarlo/Dropbox/repo/nd_language_ai_project/src/nd_language_ai/formal/computer/programming/tex/latex/part/control_sequence/kind/semtag/semtag.tex}

% for links
\input{/home/donkarlo/Dropbox/repo/nd_language_ai_project/src/nd_language_ai/formal/computer/programming/tex/latex/code_clips/seehere.tex}

\title{Particle fileter}
\date{}
\author{Mohammad Rahmani}

\begin{document}
    \maketitle
    \cite{arulampalam-2002-a-tutorial-on-particle-filters-for-online-nonlinear-non-gaussian-bayesian-tracking} describes the particle filter as a sequential Monte Carlo method that represents a probability distribution by a finite set of weighted samples called particles. The particles are propagated through the state-transition model and their weights are updated according to new observations, allowing the posterior state distribution to be estimated recursively. Resampling concentrates the particles in high-probability regions, making particle filters particularly suitable for nonlinear and non-Gaussian state-estimation problems.
    
    \cite{kutschireiter-2017-nonlinear-bayesian-filtering-and-learning-a-neuronal-dynamics-for-perception} proposes the Neural Particle Filter as a sampling-based nonlinear Bayesian filtering mechanism for perception. Particles approximate the posterior distribution and can be directly associated with the activities of neurons in a recurrent neural network. The model integrates noisy sensory information over time and across multiple sensory modalities to estimate hidden states of the environment.
    
    \section{In human body}
        Sampling-based Bayesian inference in recurrent circuits of stochastic spiking neurons \seeherefooter{https://www.nature.com/articles/s41467-023-41743-3?utm_source=chatgpt.com}
        \bibliographystyle{plainnat}
        \bibliography{/home/donkarlo/Dropbox/repo/nd_shared_research/bibliography/bibliography.bib}
\end{document}